\section{Data, conventions, and the precise conclusions}\label{sec:statements}

Throughout the paper $n\ge4$, $k=n-1$, and
\begin{equation}\label{eq:dimension-constants}
 \omega=\Area(\mathbb S^k,\sigma_k),\qquad
 c_n=\sqrt{\frac{k(k-1)}2}.
\end{equation}
The symmetrization of a covariant two-tensor includes a factor $1/2$.
We use $\Delta=\Div\nabla$, whose principal symbol is negative.
All boundary normals in expansions point into the exterior, toward its distinguished end.

\subsection{Initial data and electromagnetic normalization}

\begin{definition}[Admissible purely electric exterior]\label{def:data}
Let $\Omega$ be a connected oriented smooth $n$-manifold with nonempty compact smooth boundary $S$. The metric $g$, symmetric covariant two-tensor $K$, and vector field $E$ are smooth up to $S$. The metric space $(\Omega,g)$, including $S$, is complete. Outside a compact set there is exactly one coordinate end
\[
 \Omega\setminus C\simeq\{x\in\R^n:|x|>R_0\}.
\]
In particular the remaining topology is compact; there are no additional ends.
Put $\tau=\tr_gK$ and define
\begin{equation}\label{eq:matter-constraints}
 \begin{split}
 2\mu_m&=R_g+\tau^2-|K|_g^2-k(k-1)|E|_g^2,\\
 (J_m)_i&=\nabla^j(K_{ij}-\tau g_{ij}).
 \end{split}
\end{equation}
We require
\begin{equation}\label{eq:charged-dec}
 \mu_m\ge |J_m|_g,\qquad \Div_gE=0,\qquad
 \mu_m,\ |J_m|_g\in L^1(\Omega,\dd V_g).
\end{equation}
For some $q>(n-2)/2$, with $\rho=|x|$, assume
\begin{equation}\label{eq:weak-asymptotics}
 g_{ij}-\delta_{ij}=O_2(\rho^{-q}),\qquad
 K_{ij}=O_1(\rho^{-1-q}),\qquad
 E^i=O_1(\rho^{-k}).
\end{equation}
Here $O_j(\rho^{-a})$ bounds every coordinate derivative of order $\ell\le j$ by $C_\ell\rho^{-a-\ell}$; it places no restriction on higher derivatives beyond smoothness.
The following finite limits are assumed to exist:
\begin{align}
 \mathcal E_{\mathrm{ADM}}
 &=\frac1{2k\omega}\lim_{R\to\infty}
   \int_{\rho=R}(\partial_jg_{ij}-\partial_ig_{jj})
               (n_\delta)^i\,\dd A_\delta,\label{eq:adm-energy}\\
 (\mathbf P_{\mathrm{ADM}})_i
 &=\frac1{k\omega}\lim_{R\to\infty}
   \int_{\rho=R}(K_{ij}-\tau g_{ij})(n_\delta)^j\,\dd A_\delta,
   \label{eq:adm-momentum}\\
 Q&=\frac1\omega\lim_{R\to\infty}
   \int_{\rho=R}g(E,\nu_R)\,\dd A_g.\label{eq:charge}
\end{align}
The ADM fluxes use Euclidean coordinate normals and area, whereas the charge flux uses $g$. No timelikeness assumption is imposed on the ADM vector.
On each connected component $S_a$ choose a fixed sign $\epsilon_a\in\{1,-1\}$ and require
\begin{equation}\label{eq:mixed-trapping}
 H_{S_a}+\epsilon_a\tr_{S_a}K\le0,
 \qquad H_{S_a}=\Div_{S_a}\nu.
\end{equation}
No condition is placed on the other expansion. The signs need not agree between components.
\end{definition}

The electric field in Definition~\ref{def:data} is normalized. Its decay makes $|E|^2$ integrable, so the total constraint energy
\begin{equation}\label{eq:total-constraints}
 \mu=\mu_m+c_n^2|E|^2=\frac12(R_g+\tau^2-|K|^2),
 \qquad J=J_m,
\end{equation}
is integrable as well. The closed flux form
\begin{equation}\label{eq:electric-flux-form}
 \alpha=i_E\dd V_g
\end{equation}
is convenient when changing metrics. No source-free extension of $\alpha$ across $S$ is assumed.

\begin{definition}[Spacetime conventions]\label{def:normalization}
A Lorentz metric $G$ has signature $(-,+,\ldots,+)$. Its electromagnetic two-form $\mathcal F$ is normalized by
\begin{equation}\label{eq:einstein-maxwell}
 \Ric_G-\frac{R_G}2G
 =2\left(\mathcal F_{ac}\mathcal F_b{}^c
           -\frac14\mathcal F_{cd}\mathcal F^{cd}G_{ab}\right)+T_m,
 \qquad \dd\mathcal F=\dd(*_G\mathcal F)=0.
\end{equation}
For a spacelike embedding $\iota$ with future unit normal $N$, the induced conventions are
\begin{equation}\label{eq:induced-data}
 K(U,V)=G(\nabla_U N,V),\qquad
 E^\flat=c_n^{-1}\iota^*(i_N\mathcal F),\qquad
 \iota^*\mathcal F=0.
\end{equation}
The induced orientation is $\iota^*(i_N\dd V_G)$, and $T_m(N,N)=\mu_m$, with momentum constraint $J_m$ as in \eqref{eq:matter-constraints}. Thus the last condition in \eqref{eq:induced-data} is the purely electric condition on the initial slice. In dimensions above three, $E$ is not the unrescaled Faraday contraction.
\end{definition}

\subsection{Full enclosing cuts}

\begin{definition}[Full cut and enclosing area]\label{def:full-cut}
A full enclosing cut is $\Gamma=\partial D$, the entire intrinsic manifold boundary of a connected smooth codimension-zero submanifold-with-boundary $D\subset\Omega$ which is closed in $\Omega$, has manifold interior in $\operatorname{int}\Omega$, and contains the whole sufficiently distant coordinate end. The cut is compact, smooth, embedded and two-sided. Every component and every portion coincident with $S$ is counted. In particular $D=\Omega$ contributes the cut $S$.
Orient each component toward the end and set
\begin{equation}\label{eq:enclosing-area}
 a=A_{\min}=\inf_\Gamma\Area_g(\Gamma),\qquad
 r_A=(a/\omega)^{1/k},\qquad X_A=r_A^{k-1}.
\end{equation}
The infimum refers to the original metric $g$. Neither its positivity nor the existence of a smooth minimizer is part of this definition.
\end{definition}

Stokes' theorem for the region between a full cut and a large coordinate sphere gives
\begin{equation}\label{eq:full-cut-flux}
 \int_\Gamma\alpha=\omega Q.
\end{equation}
Component fluxes may have either sign. The adjective ``full'' is essential: replacing $\partial D$ by an empty relative frontier would discard the obstacle and change the problem.

\subsection{The precise neutral input}

We record exactly the numerical result used below. Its stronger asymptotic assumptions, sign convention and strict timelikeness are retained at each application.

\begin{theorem}[Neutral numerical input]\label{thm:neutral-input}
Let $d\ge3$ and let $(U^d,h,L)$ be a connected oriented smooth initial-data exterior with nonempty compact smooth boundary $T$, complete with $T$ included, with compact complement of exactly one Euclidean coordinate end. Write
\[
 2\mu=R_h+(\tr_hL)^2-|L|_h^2,
 \qquad J=\Div_h(L-(\tr_hL)h).
\]
Suppose $\mu\ge|J|_h$, $\mu,|J|_h\in L^1$, and for some
\[
 \frac{d-2}{2}<q_0<d-2
\]
assume
\[
 h-\delta=O_6(r^{-q_0}),\qquad L=O_5(r^{-1-q_0}).
\]
Assume that the ADM limits exist, with energy factor
$[2(d-1)\omega_{d-1}]^{-1}$ and momentum factor
$[(d-1)\omega_{d-1}]^{-1}$ applied to the Euclidean fluxes in
\eqref{eq:adm-energy}--\eqref{eq:adm-momentum}, with $g,K,\tau$ replaced by $h,L,\tr_hL$.
Assume $\mathcal E_{\mathrm{ADM}}>|\mathbf P_{\mathrm{ADM}}|$, and suppose every component of $T$ satisfies
\[
 H_T+\tr_TL\le0
\]
with normal toward the end. Let $A_*>0$ be the infimum of the areas of full cuts in the sense of Definition~\ref{def:full-cut}, now in $(U,h)$, with all components and contact with $T$ counted. Then
\begin{equation}\label{eq:neutral-input}
 \sqrt{\mathcal E_{\mathrm{ADM}}^2-|\mathbf P_{\mathrm{ADM}}|^2}
 \ge\frac12\left(\frac{A_*}{\omega_{d-1}}\right)^{(d-2)/(d-1)}.
\end{equation}
No minimizing cut is required to be smooth or attained. No spin, symmetry, maximality, vacuum or stationary-development hypothesis is imposed.
\end{theorem}

Theorem~\ref{thm:neutral-input} is the numerical part of the neutral companion result \cite[Definitions 1.1--1.2 and Theorem 1.3]{NeutralV6S74}. It is an input to this paper. Its charged extension, its weak-decay extension, its mixed-sign extension, the exclusion of a null endpoint, and every elliptic construction below are proved here. No auxiliary lemma from either companion paper is implicitly included in this input. In particular, the three-dimensional charged theorem \cite{ChargedV6S818} is not used as a premise.

\subsection{Numerical theorem and original-data equality}

\begin{theorem}[Purely electric spacetime Penrose inequality]\label{thm:numerical}
Assume Theorem~\ref{thm:neutral-input}. For every exterior in Definition~\ref{def:data}, with full enclosing area as in Definition~\ref{def:full-cut},
\begin{equation}\label{eq:numerical-conclusions}
 0<a\le\Area_g(S),\qquad
 \mathcal E_{\mathrm{ADM}}>|\mathbf P_{\mathrm{ADM}}|.
\end{equation}
Consequently the invariant mass
\[
 m=\sqrt{\mathcal E_{\mathrm{ADM}}^2-|\mathbf P_{\mathrm{ADM}}|^2}
\]
is well defined and positive, and
\begin{equation}\label{eq:charged-bound}
 m\ge |Q|,\qquad X_A\le m+\sqrt{m^2-Q^2}.
\end{equation}
Equivalently, the mass-charge bound applies when $0<X_A\le|Q|$, whereas
\begin{equation}\label{eq:polynomial-branch}
 m\ge\frac12\left(X_A+\frac{Q^2}{X_A}\right)
 \qquad\text{when }X_A>|Q|.
\end{equation}
For $Q=0$ this is $m\ge X_A/2$.
\end{theorem}

\begin{definition}[Connected horizon class]\label{def:rigidity-class}
An exterior in Definition~\ref{def:data} is in the connected horizon class if $S$ is connected,
\begin{equation}\label{eq:horizon-hypotheses}
 H_S+\tr_SK=0,\qquad
 \Area_g(\Gamma)\ge A:=\Area_g(S)>0
 \quad\text{for every full cut }\Gamma,
\end{equation}
and there is no smooth compact embedded two-sided full enclosing hypersurface entirely in $\operatorname{int}\Omega$, possibly disconnected, whose future outward expansion is nonpositive everywhere. No strictly positive stability eigenvalue, spherical topology, simple connectivity or stationary development is assumed.
\end{definition}

For $M>|q_0|$, put
\begin{equation}\label{eq:rnt-model}
 \begin{split}
 f_{M,q_0}(r)&=1-\frac{2M}{r^{n-2}}+\frac{q_0^2}{r^{2n-4}},\\
 G_{M,q_0}&=-f_{M,q_0}\dd t^2+f_{M,q_0}^{-1}\dd r^2+r^2\sigma_k,\\
 \mathcal F_{M,q_0}&=c_nq_0r^{-k}\dd t\wedge\dd r,\\
 r_+(M,q_0)&=\left(M+\sqrt{M^2-q_0^2}\right)^{1/(n-2)}.
 \end{split}
\end{equation}
The static coordinates describe $r>r_+$. The future horizon is understood in the regular ingoing or Kruskal extension, with orientation
\[
 \dd V_G=r^k\dd t\wedge\dd r\wedge\dd A_{\sigma_k}.
\]
On a static slice the normalized field is
$E=q_0r^{-k}\sqrt{f_{M,q_0}}\,\partial_r$.

\begin{theorem}[Original-data rigidity]\label{thm:rigidity}
Assume Theorem~\ref{thm:neutral-input}. Let the original data lie in the connected horizon class, suppose $m>|Q|$, and assume equality on the outer branch:
\begin{equation}\label{eq:rigidity-equality}
 \left(\frac A\omega\right)^{(n-2)/(n-1)}
 =m+\sqrt{m^2-Q^2}.
\end{equation}
Then the entire original $(\Omega,g,K,E)$, including $S$, is induced by a global smooth orientation-preserving spacelike embedding into one Reissner--Nordstr\"om--Tangherlini exterior \eqref{eq:rnt-model} with parameters $(m,Q)$ and its regular future-horizon extension. The embedding induces precisely the conventions in Definition~\ref{def:normalization}; in particular its induced magnetic two-form vanishes. Moreover $\mu_m=J_m=0$ throughout $\Omega$.

The image lies in the domain of outer communication together with the corresponding future horizon. The boundary maps either to a full smooth section of that horizon, meeting each generator once, or to the bifurcation sphere. The embedding and its future unit normal extend smoothly to $S$, and the distinguished end approaches the corresponding spatial infinity. No assertion is made behind $S$, in the extremal case, or for disconnected equality boundaries.
\end{theorem}

\begin{theorem}[Converse and sharp examples]\label{thm:converse}
Let $M>|q_0|$ and let a smooth spacelike exterior hypersurface of \eqref{eq:rnt-model} have boundary a full section of the corresponding future horizon or its bifurcation sphere. Suppose its induced normalized initial data satisfy all the hypotheses of Definitions~\ref{def:data} and~\ref{def:rigidity-class}, including the vanishing induced magnetic two-form, and suppose its actual invariant ADM mass and signed electric flux are $M$ and $q_0$. Then
\begin{equation}\label{eq:converse-area}
 A_{\min}=\Area_g(S)
 =\omega\left(M+\sqrt{M^2-q_0^2}\right)^{(n-1)/(n-2)},
\end{equation}
and equality holds in Theorem~\ref{thm:numerical}. For every $n\ge4$ and every $M>|q_0|$, static examples and admissible examples with nonzero second fundamental form exist.
\end{theorem}

\begin{remark}[Purely electric slices]
The numerical theorem allows arbitrary ADM momentum. The equality classification includes precisely the model slices satisfying its field constraints. A generic boosted slice of a charged spacetime has a nonzero induced magnetic two-form. The examples in Theorem~\ref{thm:converse} therefore use compact radial time perturbations of a static slice; their purely electric character is checked directly.
\end{remark}

\subsection{Organization and order of implication}

The proof proceeds in the following order. The exact neutral premise is Theorem~\ref{thm:neutral-input}.
\begin{center}
\small
\begin{tabular}{@{}p{.22\linewidth}p{.58\linewidth}p{.13\linewidth}@{}}
\toprule
Stage & Result used by the next stage & Sections\\
\midrule
Numerical deformation & Positive area and timelike ADM vector; charged bound of Theorem~\ref{thm:numerical} & \ref{sec:preparation}--\ref{sec:flux-conversion}\\[4pt]
Original-data variation & Lapse, shift and potential; elimination of the interior area multiplier & \ref{var:section}--\ref{var:singular-multipliers}\\[4pt]
Static reduction & Positive Killing norm and a complete vacuum static base with regular boundary & \ref{stat:section}\\[4pt]
Vacuum uniqueness & Global round identification; zero-mass rigidity uses the completed numerical theorem & \ref{uniq:section}\\[4pt]
Original-data recovery & Theorem~\ref{thm:rigidity} and the sharp examples of Theorem~\ref{thm:converse} & \ref{rec:section}\\
\bottomrule
\end{tabular}
\end{center}
The stationary fields are constructed without assuming a stationary development. The proof of zero-mass rigidity for the doubled vacuum metric uses only the numerical theorem already established at that point, so this later use introduces no circular dependence on Theorem~\ref{thm:rigidity}.
